\chapter{Cylindres nonadiques, identité et noyau de Dirac}

\section{Système inverse de cylindres nonadiques compatibles avec mémoire généalogique}

Soit
\[
X=(\Znine)^{\mathbb N}
\]
avec la topologie produit et la mesure uniforme \(\mu\). Pour
\(x=(x_1,x_2,\ldots)\), le cylindre de profondeur \(n\) est
\[
C_n(x)=\{y\in X:y_j=x_j\text{ pour }1\le j\le n\}.
\]
Les cylindres sont simultanément ouverts et fermés, forment une base et
satisfont
\[
\mu(C_n(x))=9^{-n},\qquad
C_{n+1}(x)\subset C_n(x),\qquad
\bigcap_n C_n(x)=\{x\}.
\]
Ajouter un chiffre ne crée pas un nouveau nombre : cela restreint le cylindre des
prolongements admissibles. Cette assertion donne un contenu topologique à la thèse
généalogique du nombre.

\section{Concentration cylindrique de l'identité et construction du noyau de Dirac}

\begin{theorem}[Concentration unitaire nonadique]
Définissons
\[
f_{n,x}=9^n\mathbf1_{C_n(x)}.
\]
Alors
\[
\mu(C_n(x))\to0,\qquad f_{n,x}(x)\to\infty,\qquad
\int_Xf_{n,x}\,\dd\mu=1,
\]
et les mesures \(f_{n,x}\mu\) convergent faiblement-* vers \(\delta_x\).
\end{theorem}

\begin{proof}
La première identité est immédiate. La hauteur de \(f_{n,x}\) sur son support
est \(9^n\), de sorte qu'elle diverge. L'intégrale vaut
\(9^n9^{-n}=1\). Pour \(g\in C(X)\), l'oscillation de \(g\) dans \(C_n(x)\)
tend vers zéro par continuité uniforme ; par conséquent
\[
\left|\int g f_{n,x}\,\dd\mu-g(x)\right|
\le \sup_{y\in C_n(x)}|g(y)-g(x)|\longrightarrow0.
\]
C'est la convergence faible-* vers \(\delta_x\).
\end{proof}

Le théorème n'affirme pas \(1/0=\infty\). Il exhibe trois évaluations d'un même
objet : mesure du support, densité et masse totale. Par blocs de trois chiffres,
\[
\mu(C_{3m})=729^{-m},\qquad
f_{3m}=729^m\mathbf1_{C_{3m}},\qquad
729^{-m}729^m=1.
\]
La conservation évolutive de l'unité prend ici une forme analytique exacte.

\section{Convergence du noyau de Kronecker vers la distribution de Dirac}

Au niveau fini \(X_n=(\Znine)^n\), soit
\[
K_n(x,y)=9^n\mathbf1_{\{x=y\}}.
\]
Avec la mesure uniforme, l'opérateur
\[
(E_nf)(x)=\sum_{y\in X_n}K_n(x,y)f(y)9^{-n}
\]
est l'identité. Dans la limite cylindrique, la condition \(x=y\) est remplacée par
« même préfixe de longueur \(n\) » et \(E_n\) est l'espérance conditionnelle sur
l'algèbre cylindrique. On a
\[
E_n^2=E_n,\qquad E_n^*=E_n,\qquad E_nf\to f
\quad\text{dans }L^2(X,\mu).
\]

Le pont exact est donc
\[
\delta_{ij}
\longrightarrow K_n(x,y)
\longrightarrow\delta_x.
\]
Le delta de Kronecker est un noyau matriciel ; le delta de Dirac est une mesure ou
une distribution. Aucun n'est l'opérateur de Dirac
\(D=\gamma^\mu\partial_\mu\). Le passage aux spineurs et à Clifford exige
un second opérateur d'entrelacement et sera traité ultérieurement.

\section{Loi d'échelle, dimension effective et densité de la mesure cylindrique}

Supposons qu'une contraction de rapport \(\lambda\) agisse de manière indépendante
dans \(d\) directions et que la mesure produit varie comme \(\lambda^d\). Pour
conserver une masse égale à un, la densité doit varier comme
\[
\rho\longmapsto\lambda^{-d}\rho.
\]
Les exposants \((-1,-2,-3)\) apparaissent alors comme des poids de densité en ligne,
en aire et en volume. Cette proposition est générale et exacte sous l'hypothèse de
produit. Elle ne permet pas d'identifier automatiquement \(\lambda\) à
\(\alpha_{\rm HMT}\) ; cette identification nécessite une carte métrique.

\section{Monodromie nonadique et calculabilité de l'évaluation archimédienne}

L'existence d'une branche infinie dans un arbre à branchement fini peut
découler du lemme de König, mais König ne produit pas de sélecteur calculable. Il existe
des arbres binaires récursifs infinis sans branche récursive. La connexion nonadique
conjointe ne repose pas sur cette existence abstraite : elle possède un sélecteur uniforme
certifié qui engendre à profondeur finie arbitraire et n'ouvre les témoins
décimaux qu'après la génération.

\begin{falsadorBox}[title={Trois confusions interdites}]
On n'écrira pas \(1/0\) ; on ne représentera pas le delta de Dirac comme une fonction
ordinaire avec « un un infini » ; on n'utilisera pas König comme algorithme. Le résultat
fort est la compatibilité exacte entre cylindres, densité unitaire et limite
faible, ainsi que le sélecteur constructif indépendant de la comparaison décimale.
\end{falsadorBox}
